\chapter[Marco de agentes]{Diseño del marco de trabajo basado en agentes}
\label{cap:agentes}

El marco organiza una búsqueda experimental sobre un espacio de soluciones definido para cada tarea. Recibe un conjunto de datos con particiones fijadas, métricas de calidad y recursos, restricciones sobre las salidas y un presupuesto. Produce candidatos ejecutables, resultados verificables y un historial que permite reconstruir por qué se eligió cada experimento. La arquitectura se inspira en el ciclo de propuesta, ejecución y análisis de ASI-Arch \citep{liu2025asiarch}, pero cambia el objeto de búsqueda: aquí se comparan soluciones para detección, \emph{spoiling} y su integración, no mecanismos de atención lineal.

\section{Interfaz de tarea}

La unidad de adaptación es una especificación de tarea $T=(D,S,M,V,B)$, donde $D$ identifica datos y particiones, $S$ delimita el espacio de soluciones, $M$ define las métricas de selección y reporte, $V$ reúne verificaciones y restricciones, y $B$ fija el presupuesto. La especificación incluye formato de entrada y salida, identificadores de ejemplos, semillas, tiempo máximo y recursos permitidos. Los agentes consultan esta interfaz; la lógica del ciclo no necesita conocer si una salida es una clase o una respuesta textual.

En detección, $M$ incluye F1 de la clase clickbait y latencia de inferencia en un entorno fijado. En \emph{spoiling}, incluye métricas automáticas y costo de las consultas; la fidelidad a la fuente y la utilidad de la respuesta requieren además un control humano. El espacio $S$ puede incluir líneas base léxicas, codificadores ajustados, recuperación y generación. Las variantes con imagen u otras modalidades solo se comparan cuando las entradas y las etiquetas están disponibles. Cada tarea conserva sus restricciones, por ejemplo la longitud máxima de una respuesta.

\section{Ciclo experimental}

En cada iteración, el agente proponente recibe la especificación de tarea y un resumen del historial. Devuelve una hipótesis concreta, un candidato o cambio de configuración, su relación con intentos anteriores y una estimación de costo. La propuesta solo avanza si respeta el espacio y el presupuesto definidos. Las decisiones de un modelo de lenguaje no sustituyen verificaciones deterministas sobre datos, código y recursos.

El ejecutor prepara el entorno, comprueba el formato de entradas y salidas, lanza el entrenamiento o la inferencia y evalúa sobre la partición autorizada. Debe registrar tanto corridas válidas como errores, tiempos agotados y restricciones incumplidas. El analista compara resultados bajo el mismo protocolo y devuelve al proponente un resumen de mejoras, fallos y líneas de búsqueda descartadas. La memoria de experimentos conserva identificador, configuración, versión del código y de los datos, semillas, métricas, consumo de recursos y artefactos de predicción.

La condición de parada combina el presupuesto total con límites por candidato. El resultado de una búsqueda es el mejor candidato \emph{válido} según el criterio fijado antes de ejecutarla, junto con la trayectoria de búsqueda y su costo. No se selecciona el resultado por inspección del conjunto final de prueba.

La búsqueda puede considerar configuraciones del pipeline completo: umbral del detector, decisión de cuándo ofrecer un \emph{spoiler}, recuperación de la fuente y generador. Para interpretar los resultados se miden por separado detector y generador antes de evaluar su combinación. Las salidas se vinculan al identificador de la publicación y a la versión del contenido recuperado; una respuesta basada en un artículo que cambió o no se pudo obtener no se trata como una evaluación normal del mismo ejemplo.

\section{Adquisición de datos multimodales}

Se propone un subsistema de agentes de recolección para construir una extensión de TA1C con publicaciones de distintas fuentes y medios. Su función es localizar y recuperar publicaciones, noticias y archivos asociados, extraer texto y metadatos, y registrar procedencia y fallos. Los conectores de cada fuente deben producir un esquema común, con identificadores, URL, fecha de acceso, tipo de medio, relación con la publicación, estado de recuperación y huella del contenido. La extracción automática no asigna por sí sola etiquetas de clickbait ni respuestas de referencia: esas decisiones requieren un protocolo de anotación y control de calidad.

Este subsistema es distinto del ciclo que busca soluciones. Los datos recolectados se congelan y auditan antes de que el proponente reciba una partición de desarrollo. La disponibilidad, las condiciones de uso y el derecho a conservar o distribuir medios pueden limitar el corpus resultante; por eso el tamaño y la cobertura de la extensión no se presuponen aquí.

\section{Controles de validez y reproducibilidad}

Las particiones de evaluación quedan fuera del contexto de los agentes durante la búsqueda. Los candidatos se validan con comprobaciones de esquema, ausencia de identificadores prohibidos, límites de recursos y reglas de la tarea. Las ejecuciones fallidas cuentan para el presupuesto: excluirlas favorecería indebidamente a un método que propone más soluciones inválidas. El registro debe permitir repetir un candidato sin depender del estado de una conversación o de una sesión interactiva.

La comparación experimental separa el efecto de proponer candidatos del de ejecutar más pruebas. Por ello se igualan los recursos disponibles a la búsqueda aleatoria y se realizan ablaciones de componentes del ciclo: sin análisis histórico, sin memoria recuperable o con propuestas aleatorias dentro del mismo espacio. Los resultados del estudio de base de los capítulos siguientes sirven para definir entradas, restricciones y referencias; no se consideran una evaluación del marco.

% Para cerrar este capítulo con la implementación real: documentar componentes,
% modelos y versiones de agentes, prompts, esquemas de datos, aislamiento de
% ejecución, política de errores, repositorio y diagrama de arquitectura.
